\section{Additional Experimental Details}
\paragraph{Base setting.}
All experiments use an RWKV7 backbone with prompt-time Future-Seed only: generation remains left-to-right and causal at decode time. The real-task campaign uses a single NVIDIA RTX 4090 and serial quick-to-medium screening to conserve wall-clock budget.

\paragraph{Post-training search policy.}
The late-stage queues use a two-stage policy: a quick stage for baseline and candidate screening, and a medium stage only for candidates that clear a positive gate. In the final fastdiscover queues, we prune candidates below $-0.5$pp at quick stage and only promote a top candidate above $+0.8$pp to medium budget.

\paragraph{Task families.}
The real-task audit covers protein secondary-structure labeling, multi-hop QA, code-generation-style long-context completion, ARC multiple choice, punctuation restoration, and a small set of constraint and planning proxies. The most stable family under this setup is protein secondary-structure labeling.

\section{Additional Synthetic Results}
\begin{table*}[t]
\centering
\small
\caption{Transformer baselines on KVSORT. Hungarian decoding restores validity but not ordering.}
\label{tab:appendix-transformer}
\begin{tabular}{lcccc}
\toprule
Model & Decode & Exact & Key-valid & Key-order \\
\midrule
Transformer-MLM & argmax & 0.00\% & 0.00\% & 0.00\% \\
Transformer-MLM & hungarian & 0.00\% & 100.00\% & 0.00\% \\
Transformer-Causal & hungarian & 0.00\% & 100.00\% & 0.00\% \\
Transformer-Causal + ATTN\_FS & hungarian & 0.00\% & 100.00\% & 0.00\% \\
Transformer-MLM + Sinkhorn & hungarian & 0.00\% & 100.00\% & 0.00\% \\
\bottomrule
\end{tabular}
\end{table*}

\begin{table}[t]
\centering
\small
\caption{4x4 Sudoku solve-rate phase curve. Future-Seed delays the collapse point to harder puzzles.}
\label{tab:appendix-sudoku-phase}
\begin{tabular}{rcc}
\toprule
Holes & FS=0 solve & FS=1 solve \\
\midrule
4 & 35.30\% & 100.00\% \\
6 & 16.65\% & 100.00\% \\
8 & 5.20\% & 100.00\% \\
10 & 0.95\% & 95.10\% \\
12 & 0.00\% & 55.10\% \\
14 & 0.00\% & 0.00\% \\
\bottomrule
\end{tabular}
\end{table}

\begin{table}[t]
\centering
\small
\caption{A simple Sudoku consistency regularizer helps at holes=12 but does not solve the hardest setting.}
\label{tab:appendix-sudoku-consistency}
\begin{tabular}{lccc}
\toprule
Setting & holes=10 & holes=12 & holes=14 \\
\midrule
lambda=0.0 & 93.85\% & 47.95\% & 0.00\% \\
lambda=0.1 & 95.15\% & 58.55\% & 0.00\% \\
\bottomrule
\end{tabular}
\end{table}


\paragraph{Takeaway.}
The synthetic suite supports a narrow but strong claim: Future-Seed is especially helpful when a recurrent model must repair a masked region using information that arrives on the right and must also satisfy a global consistency constraint. The Transformer baselines in Table~\ref{tab:appendix-transformer} show that simply restoring bidirectional visibility or decoding with Hungarian matching is not enough on KVSORT; the useful signal comes from the specific cross-layer state feedback path.

\section{Additional Post-Training Results}
\begin{table}[t]
\centering
\small
\caption{Additional task families discussed in appendix.}
\label{tab:appendix-negative-tasks}
\begin{tabular}{lcccc}
\toprule
Task family & Best gain & Median gain & Positive med count & Judgment \\
\midrule
graph\_color & +8.33pp & +0.00pp & 4/10 & diagnostic only \\
tsp\_mask & +25.00pp & +2.08pp & 2/4 & appendix only \\
countdown & -4.17pp & -4.17pp & 0/2 & negative \\
nqueens & -4.17pp & -10.42pp & 0/2 & negative \\
sat3 & +0.00pp & +0.00pp & 0/4 & near-zero \\
\bottomrule
\end{tabular}
\end{table}


\paragraph{Closure window.}
The final narrow search focused on \texttt{mbpp\_longctx} and \texttt{arc\_mc} before a breadth pivot. The best closure-window results were \texttt{mbpp\_longctx\_seed38\_depthmix\_anchor} at $+5.01$pp and \texttt{arc\_mc\_seed260\_depthmix\_qfirst} at $+9.38$pp, but neither family held up as a clean held-out confirmation line. The strict confirmation queue failed on \texttt{mbpp\_longctx} and flipped negative on \texttt{arc\_mc}.

\paragraph{Breadth pivot.}
Rounds 805--808 were used to stop over-spending on the same \texttt{mbpp\_longctx} recipe family and test a wider set of real tasks. Only \texttt{hotpot} converted to medium-stage positives in that pivot, with a best gain of $+1.00$pp. \texttt{protein\_ss} and \texttt{squad} stayed near the promote gate but did not convert in those breadth rounds.

\section{Reproducibility and Artifact Layout}
The supplementary material is intended to include:
\begin{itemize}
    \item the code for the synthetic task suite under \texttt{rwkv-diff-future-seed/};
    \item the post-training launchers under \texttt{posttrain\_rwkv7/scripts/};
    \item the final queue files used in the closure and breadth pivot;
    \item the synced summaries and JSONL records for rounds 783--808 under \texttt{posttrain\_rwkv7/runs/};
    \item this paper package under \texttt{paper/neurips2025/}.
\end{itemize}
A compact build path is provided in the paper package via \texttt{build.sh}. The tables used in the paper are regenerated from committed paper-side metrics using \texttt{render\_tables.py}.

\section{Assets, Licenses, and Scope}
The real-task audit uses publicly available datasets and benchmark snapshots, including ARC, HotpotQA, SQuAD, MBPP, and a public protein secondary-structure dataset snapshot. We cite the original benchmark creators in the main paper. We do not claim to have completed a full legal audit of every upstream redistribution license inside this submission; this is one reason the checklist answer on explicit license coverage is conservative.
